% Editorial master: guides/pronunciation/student.md

% Master SHA-256: 7c6b8e6907962362381148b077dbbbb23b6fbe4196f98064faef00d2988ec346

% Initial assembly; retain complete reviewed TeX for reproduction.

\newcommand{\runningtitle}{Pronunciation}

% Layout inherited from the reviewed previous course; set a running title per unit.
% This is a small hand-maintained TeX layout, not a conversion pipeline.
\RequirePackage{fix-cm}
\documentclass[11pt,twoside]{article}
\usepackage{fontspec}
\setmainfont{Cambria}[
  Renderer=HarfBuzz,
  BoldFont=cambriab.ttf,
  ItalicFont=cambriai.ttf,
  BoldItalicFont=cambriaz.ttf
]
\setmonofont{lmmono10-regular.otf}[BoldFont=lmmonolt10-bold.otf]
\usepackage[paperwidth=165mm,paperheight=241mm,inner=21mm,outer=19mm,
  top=22mm,bottom=20mm,headheight=14pt,headsep=7mm,footskip=10mm]{geometry}
\usepackage{xcolor,array,booktabs,longtable,tabularx}
\usepackage{enumitem,fancyhdr}
\usepackage[unicode,hidelinks]{hyperref}
\definecolor{ink}{HTML}{000000}
\definecolor{muted}{HTML}{4C5558}
\setlength{\parindent}{0pt}
\setlength{\parskip}{3.5pt}
\AtBeginDocument{\fontsize{11}{14.3}\selectfont}
\setlength{\tabcolsep}{4pt}
\renewcommand{\arraystretch}{1.10}
\setlist[itemize]{leftmargin=1.3em,label=\textbullet,itemsep=2pt,topsep=3pt,parsep=0pt}
\setlist[enumerate]{leftmargin=1.8em,itemsep=3pt,topsep=4pt,parsep=0pt}
\widowpenalty=10000
\clubpenalty=10000
\emergencystretch=1em
\raggedbottom
\pagestyle{fancy}
\fancyhf{}
\providecommand{\runningtitle}{Old English}
\fancyhead[LE,RO]{\small\itshape \runningtitle}
\fancyhead[LO,RE]{\small Old English}
\fancyfoot[LE,RO]{\small\thepage}
\renewcommand{\headrulewidth}{0.3pt}
\newcommand{\booktitle}[2]{
  {\fontsize{9}{11}\selectfont\color{muted}#1\par}
  \vspace{3pt}
  {\fontsize{20}{24}\selectfont\color{ink}#2\par}
  \vspace{6pt}
}
\newcommand{\source}[1]{\par{\footnotesize\color{muted}#1\par}}
\newcommand{\readingtext}[1]{
  \par\addvspace{5pt}
  \begingroup\leftskip=1.5em
    \fontsize{12}{16}\selectfont #1\par
  \endgroup\addvspace{5pt}
}
\makeatletter
\renewcommand{\section}{\@startsection{section}{1}{\z@}
  {-2.4ex \@plus -0.6ex \@minus -0.2ex}{0.8ex \@plus 0.2ex}
  {\normalfont\fontsize{12}{15}\selectfont\bfseries\color{ink}}}
\renewcommand{\subsection}{\@startsection{subsection}{2}{\z@}
  {-1.8ex \@plus -0.5ex \@minus -0.2ex}{0.5ex \@plus 0.2ex}
  {\normalfont\fontsize{11}{14}\selectfont\bfseries\color{ink}}}
\makeatother


\hyphenpenalty=10000

\exhyphenpenalty=10000

\emergencystretch=3em

\newcommand{\reserve}[1]{\par\penalty-100\begingroup\dimen0=\pagegoal\advance\dimen0 by-\pagetotal\ifdim\dimen0<#1\newpage\fi\endgroup}

\begin{document}

\addfontfeatures{Ligatures=TeXOff}

\raggedright

\booktitle{GUIDE}{A guide to Old English pronunciation}

% mdparagraph:P001:begin

This guide gives a practical pronunciation for reading the course. English examples refer to a particular sound in an ordinary American pronunciation, not to the whole word. Where English has no equivalent, the instructions describe the position of the tongue and lips. Every Old English word is glossed; its meaning need not be memorized.

% mdparagraph:P001:end



\reserve{5\baselineskip}

% mdparagraph:P002:begin

\section*{1. Letters and length}

% mdparagraph:P002:end



% mdparagraph:P003:begin

\textbf{Æ æ}, called \emph{ash}, is one vowel letter. \textbf{Þ þ}, \emph{thorn}, and \textbf{Ð ð}, \emph{eth}, represent the sounds written \textbf{th} in English \emph{thin} and \emph{this}. Either letter can represent either sound. This course normally uses \textbf{ð}. Cook's hooked \textbf{ę} and \textbf{ǫ} have the sounds of \textbf{e} and \textbf{o}; the hooks record their history, not additional sounds.[1]

% mdparagraph:P003:end



% mdparagraph:P004:begin

A \textbf{macron}, the line above a vowel, marks length: \textbf{ā, ǣ, ē, ī, ō, ū, ȳ}. Hold a long vowel appreciably longer than a short one. Long does not mean the English letter name, and a macron does not mark stress. Thus \textbf{ac}, “but”, has a short vowel; \textbf{āc}, “oak”, has a long one. Both end with the sound of \textbf{k}.

% mdparagraph:P004:end



% mdparagraph:P005:begin

In a long two-letter vowel spelling, the macron stands on the first letter: \textbf{ēa, ēo, īe}. The two letters belong to one syllable. Preserve the printed marks; a following pair of consonants does not turn an unmarked vowel into a long one.

% mdparagraph:P005:end



\reserve{5\baselineskip}

% Keep the complete simple-vowel reference table together on the next page.
\newpage

% mdparagraph:P006:begin

\section*{2. Simple vowels}

% mdparagraph:P006:end



% mdparagraph:P007:begin

Use these short and long values. Some of the practical English approximations differ in quality as well as duration.[2]

% mdparagraph:P007:end



% mdtable:T01:begin
\par\begingroup
\fontsize{10.5}{13.2}\selectfont
\setlength{\LTleft}{0pt}\setlength{\LTright}{\fill}
\setlength{\LTpre}{5pt}\setlength{\LTpost}{5pt}
\begin{longtable}{@{}>{\raggedright\arraybackslash}p{16.000mm}>{\raggedright\arraybackslash}p{51.688mm}>{\raggedright\arraybackslash}p{51.688mm}@{}}
\toprule
\textbf{Letters} & \textbf{Short vowel} & \textbf{Long vowel} \\
\midrule
\endfirsthead
\toprule
\textbf{Letters} & \textbf{Short vowel} & \textbf{Long vowel} \\
\midrule
\endhead
\bottomrule
\endfoot
\bottomrule
\endlastfoot
a, ā & The vowel of father, brief: \textbf{ac}, “but”. & Hold the same vowel: \textbf{āc}, “oak”. \\
æ, ǣ & The vowel of cat: \textbf{sæt}, “sat”. & Hold that same vowel: \textbf{sǣd}, “seed”. \\
e, ē & The vowel of bet: \textbf{męte}, “food”. & Hold the beginning of day without its final glide: \textbf{hē}, “he”. \\
i, ī & The vowel of bit: \textbf{is}, “is”. & The vowel of see, held: \textbf{īs}, “ice”. \\
o, ō & A brief rounded vowel, described below: \textbf{god}, “god”. & Hold the beginning of go without its final glide: \textbf{gōd}, “good”. \\
u, ū & The vowel of foot: \textbf{full}, “full”. & The vowel of food, held: \textbf{fūl}, “vile”. \\
y, ȳ & The tongue position of bit with rounded lips: \textbf{wyrt}, “herb; plant”. & The tongue position of see with rounded lips, held: \textbf{fȳr}, “fire”. \\
\end{longtable}
\par\endgroup
% mdtable:T01:end

% mdparagraph:P008:begin

For \textbf{o}, the vowel of \emph{caught} is a useful starting point if you distinguish \emph{cot} from \emph{caught}. If they sound alike, begin with the first part of \emph{go}, lower the jaw a little while keeping the lips rounded, and make the vowel brief. Keep \textbf{o} distinct from the unrounded \emph{father} vowel. For \textbf{ō}, keep the jaw and tongue steady: do not finish with an \textbf{oo} movement.

% mdparagraph:P008:end



% mdparagraph:P009:begin

For \textbf{ē}, likewise stop before the glide at the end of \emph{day}. Do not pronounce \textbf{ǣ} like American \emph{care}: it is the \emph{cat} vowel held longer, without an \textbf{r} sound.

% mdparagraph:P009:end



% mdparagraph:P010:begin

In a \textbf{stressed} syllable, short \textbf{a} before \textbf{m} or \textbf{n} takes the more rounded short-\textbf{o} quality just described. Use this in \textbf{mann}, “man; human being”, and \textbf{lang}, “long”. It is not the vowel of ordinary American \emph{man}. This rule does not require the same rounding in every unstressed \textbf{-an} ending.

% mdparagraph:P010:end



% mdparagraph:P011:begin

For \textbf{y} and \textbf{ȳ}, keep the tongue forward while rounding the lips as for \textbf{oo} in \emph{food}. For short \textbf{y}, begin with the tongue position of \emph{bit}; for long \textbf{ȳ}, use \emph{see}. Rounding the lips must not pull the tongue back into an ordinary \textbf{oo} vowel. These sounds have no usual American English equivalent. Here \textbf{y} is a vowel, not the consonant at the beginning of \emph{yes}.

% mdparagraph:P011:end



\reserve{5\baselineskip}

% mdparagraph:P012:begin

\section*{3. Diphthongs}

% mdparagraph:P012:end



% mdparagraph:P013:begin

A \emph{diphthong} is a vowel that moves from one position toward another within one syllable. For these spellings, make one smooth movement toward \textbf{oo}. Do not give the second letter a separate beat.[3]

% mdparagraph:P013:end



% mdtable:T02:begin
\par\begingroup
\fontsize{10.5}{13.2}\selectfont
\setlength{\LTleft}{0pt}\setlength{\LTright}{\fill}
\setlength{\LTpre}{5pt}\setlength{\LTpost}{5pt}
\begin{longtable}{@{}>{\raggedright\arraybackslash}p{17.000mm}>{\raggedright\arraybackslash}p{49.000mm}>{\raggedright\arraybackslash}p{53.377mm}@{}}
\toprule
\textbf{Spellings} & \textbf{Movement} & \textbf{Supplied examples} \\
\midrule
\endfirsthead
\toprule
\textbf{Spellings} & \textbf{Movement} & \textbf{Supplied examples} \\
\midrule
\endhead
\bottomrule
\endfoot
\bottomrule
\endlastfoot
ea, ēa & Start with the vowel of cat, then round toward \textbf{oo}. & \textbf{earm}, “arm”; \textbf{ēare}, “ear”. \\
eo, ēo & Start with the steady first vowel of day, then round toward \textbf{oo}. & \textbf{eorl}, “man”; \textbf{fēoll}, “fell”. \\
ie, īe & Start with the vowel of see, then round toward \textbf{oo}. & \textbf{giefan}, “give”; \textbf{hīeran}, “hear”. \\
\end{longtable}
\par\endgroup
% mdtable:T02:end

% mdparagraph:P014:begin

In \textbf{ēa, ēo, īe}, hold the first part longer before making the movement. The \textbf{eo} of \textbf{eorl} is short; the \textbf{ēo} of \textbf{fēoll} is long. The word \textbf{ēare} has two syllables: its long \textbf{ēa}, followed by a light final \textbf{e}. The opening \textbf{g} of \textbf{giefan} has the consonant sound of \textbf{y} in \emph{yes}.

% mdparagraph:P014:end



% mdparagraph:P015:begin

There are word-specific exceptions. In \textbf{geong}, “young”, \textbf{e} does not supply a vowel sound: the initial \textbf{g} is like \textbf{y} in \emph{yes}, followed by short rounded \textbf{o}. Its \textbf{ng} has the sound heard in \emph{finger}, including the \textbf{g}. Do not give this word the ordinary \textbf{eo} glide or copy the vowel of modern \emph{young}.

% mdparagraph:P015:end



\reserve{5\baselineskip}

% mdparagraph:P016:begin

\section*{4. Consonants}

% mdparagraph:P016:end



% mdparagraph:P017:begin

Read \textbf{b, d, l, m, n, p, t, w} as at the beginning of \emph{boy, dog, lip, man, no, pen, top, win}. Use a pronounced American \textbf{r}, as in \emph{run}, as a practical approximation. The remaining letters and combinations need the following distinctions.[4]

% mdparagraph:P017:end



\reserve{5\baselineskip}

% mdparagraph:P018:begin

\subsection*{c, g, and combinations}

% mdparagraph:P018:end



% mdtable:T03:begin
\par\begingroup
\fontsize{10.5}{13.2}\selectfont
\setlength{\LTleft}{0pt}\setlength{\LTright}{\fill}
\setlength{\LTpre}{5pt}\setlength{\LTpost}{5pt}
\begin{longtable}{@{}>{\raggedright\arraybackslash}p{23.000mm}>{\raggedright\arraybackslash}p{36.000mm}>{\raggedright\arraybackslash}p{60.377mm}@{}}
\toprule
\textbf{Spelling} & \textbf{Sound} & \textbf{Supplied examples and limits} \\
\midrule
\endfirsthead
\toprule
\textbf{Spelling} & \textbf{Sound} & \textbf{Supplied examples and limits} \\
\midrule
\endhead
\bottomrule
\endfoot
\bottomrule
\endlastfoot
c & \textbf{k} as in kit, or \textbf{ch} as in chin. & \textbf{cuman}, “come”, begins with \textbf{k}. Both \textbf{c} letters in \textbf{cyrice}, “church”, are \textbf{ch}. \\
c before a front vowel & Often \textbf{ch}, but not invariably. & \textbf{cēlan}, “cool”, begins with \textbf{k}, despite \textbf{ē}. \\
g at the beginning & Hard \textbf{g} as in go, or consonant \textbf{y} as in yes. & \textbf{god}, “god”, has hard \textbf{g}; \textbf{giefan}, “give”, begins with \textbf{y}. \\
g elsewhere & Consonant \textbf{y}, or the continuous sound described below; some contexts retain hard \textbf{g}. & \textbf{dæg}, “day”, ends with \textbf{y}; \textbf{fugol}, “bird”, has continuous \textbf{g}. \\
sc & Usually \textbf{sh} as in ship. & \textbf{scip}, “ship”, has one \textbf{sh} sound; \textbf{frosc}, “frog”, instead ends with \textbf{s} followed by \textbf{k}. \\
cg; usually gg & The final sound of bridge. & \textbf{ecg}, “edge”. The supplied exception \textbf{docga}, “dog”, has a prolonged hard \textbf{g}. \\
ng & Usually as in finger, with audible \textbf{g}. & \textbf{lang}, “long”. In \textbf{sengan}, “singe”, \textbf{ng} instead has the sound at the end of singe. \\
\end{longtable}
\par\endgroup
% mdtable:T03:end

% mdparagraph:P019:begin

The front vowels \textbf{e, i, æ, y} help explain many \textbf{ch} and \textbf{y} pronunciations, but the next written letter is not a complete rule. Words and inflected forms have exceptions. Use the supplied sound for each example; do not turn every \textbf{c} before \textbf{e} into \textbf{ch}. In \textbf{cyrice}, keep all three syllables, including the final \textbf{e}.

% mdparagraph:P019:end



% mdparagraph:P020:begin

The continuous \textbf{g} in \textbf{fugol} is unfamiliar in ordinary American English. Put the back of the tongue near the position for hard \textbf{g}, but leave a narrow gap. Let air pass continuously while the throat vibrates: the sound can be held, unlike the stopped \textbf{g} of \emph{go}. Its phonetic sign is \textbf{[ɣ]}. Do not omit it or substitute \textbf{y}.

% mdparagraph:P020:end



\reserve{5\baselineskip}

% mdparagraph:P021:begin

\subsection*{h and sounded clusters}

% mdparagraph:P021:end



% mdparagraph:P022:begin

At the beginning of a syllable, \textbf{h} is like English \textbf{h} in \emph{hat}. Later in a syllable it has a stronger, continuous sound made by narrowing the passage above the tongue.[5]

% mdparagraph:P022:end



% mdparagraph:P023:begin

In \textbf{brōhte}, “brought”, raise the back of the tongue toward the position for \textbf{k}, leaving a gap for breath. Keep the throat unvoiced. This back sound is \textbf{[x]}. Both it and the following \textbf{t} are heard.

% mdparagraph:P023:end



% mdparagraph:P024:begin

In \textbf{niht}, “night”, the preceding \textbf{i} brings the tongue farther forward. Raise the middle of the tongue toward the roof of the mouth and let breath pass without voice. This front sound is \textbf{[ç]}. Keep the \textbf{t} too: the Old English word does not have modern English's silent \textbf{gh}. Front vowels generally favour this front articulation; back vowels favour the back one.

% mdparagraph:P024:end



% mdparagraph:P025:begin

Sound the initials in \textbf{hl, hn, hr, hw, wr, wl, cn}. In \textbf{hlāf}, “loaf”, begin with audible breath before \textbf{l}; in \textbf{hwæl}, “whale”, keep \textbf{h} with \textbf{w}. In \textbf{wrītan}, “write”, pronounce \textbf{w} and \textbf{r}; in \textbf{cnīf}, “knife”, pronounce \textbf{k} and \textbf{n}. Join the consonants without putting a vowel between them. The spelling \textbf{cw} has \textbf{k} followed by \textbf{w}. Read \textbf{x} and \textbf{hs} as \textbf{k} followed by \textbf{s}; \textbf{weax}, “wax”, ends with that pair.

% mdparagraph:P025:end



\reserve{5\baselineskip}

% mdparagraph:P026:begin

\section*{5. f, s, and th: voice}

% mdparagraph:P026:end



% mdparagraph:P027:begin

Compare \emph{fan} with \emph{van}, \emph{sip} with \emph{zip}, and \emph{thin} with \emph{this}. The initial consonant of the second word has voice: the throat vibrates. Vowels and consonants such as \textbf{l, m, n, r} also carry voice.[6]

% mdparagraph:P027:end



% mdparagraph:P028:begin

The ordinary rule for single \textbf{f, s, þ/ð} is:

% mdparagraph:P028:end



% mdparagraph:P029:begin

\begin{itemize}
\item At the beginning or end of a word, or beside a voiceless consonant, use the sounds of \emph{fan, sip, thin}.
\end{itemize}

% mdparagraph:P029:end



% mdparagraph:P030:begin

\begin{itemize}
\item Inside a word, between voiced sounds, use the sounds of \emph{van, zip, this}.
\end{itemize}

% mdparagraph:P030:end



% mdtable:T04:begin
\par\begingroup
\fontsize{10.5}{13.2}\selectfont
\setlength{\LTleft}{0pt}\setlength{\LTright}{\fill}
\setlength{\LTpre}{5pt}\setlength{\LTpost}{5pt}
\begin{longtable}{@{}>{\raggedright\arraybackslash}p{20.000mm}>{\raggedright\arraybackslash}p{22.000mm}>{\raggedright\arraybackslash}p{77.377mm}@{}}
\toprule
\textbf{Word} & \textbf{Meaning} & \textbf{Sound to use} \\
\midrule
\endfirsthead
\toprule
\textbf{Word} & \textbf{Meaning} & \textbf{Sound to use} \\
\midrule
\endhead
\bottomrule
\endfoot
\bottomrule
\endlastfoot
cnīf & knife & Final \textbf{f} as in fan. \\
cnīfas & knives & \textbf{f} as in van, between vowels; final \textbf{s} as in sip. \\
rīsan & rise & \textbf{s} as in zip, between vowels. \\
hūs & house & Final \textbf{s} as in sip. \\
ōðer & other & \textbf{ð} as in this, between vowels. \\
ðā & then & Initial \textbf{ð} as in thin. \\
\end{longtable}
\par\endgroup
% mdtable:T04:end

% mdparagraph:P031:begin

Thus \textbf{is}, “is”, ends with the \textbf{s} of \emph{sip}, not the \textbf{z} of its modern English equivalent. The letter \textbf{ð} does not by itself tell you to use the voiced sound.

% mdparagraph:P031:end



% mdparagraph:P032:begin

Stress and the boundaries between word parts can preserve exceptions. In \textbf{hǣlðu}, “health”, and \textbf{forðām}, “therefore”, use the unvoiced \textbf{th} of \emph{thin}. Learn these as supplied word-specific readings; the surrounding voiced sounds do not override them.

% mdparagraph:P032:end



\reserve{5\baselineskip}

% mdparagraph:P033:begin

\section*{6. Duration, syllables, and stress}

% mdparagraph:P033:end



% mdparagraph:P034:begin

Keep vowel length separate from consonant length. In \textbf{full}, “full”, \textbf{u} is short; \textbf{fūl}, “vile”, has long \textbf{ū}. A doubled consonant at the end of a word normally has ordinary single-consonant duration: the final \textbf{ll} of \textbf{full} does not require extra holding.

% mdparagraph:P034:end



% mdparagraph:P035:begin

Between vowels, a doubled consonant has extra duration, like the \textbf{n} across the join in \emph{pen-knife}. Do not insert a vowel between its two letters.[7]

% mdparagraph:P035:end



% mdparagraph:P036:begin

With a doubled stop, hold the closure longer. In \textbf{habban}, “have”, keep the lips closed for \textbf{bb}, then release into the following vowel. There are two syllables. Do not add another vowel or another syllable to pronounce the doubling.

% mdparagraph:P036:end



% mdparagraph:P037:begin

Written unstressed vowels remain audible and distinguishable. In \textbf{męte}, “food”, there are two syllables, and final \textbf{e} is a light \emph{bet} vowel. In \textbf{fugol}, “bird”, there are also two syllables; its \textbf{o} is heard. Do not reduce every ending to the \textbf{uh} at the beginning of modern \emph{about}.

% mdparagraph:P037:end



% mdparagraph:P038:begin

The main stress normally falls on the first syllable, as in \textbf{habban}, “have”, and \textbf{heofonan}, “heaven”. Prefixes complicate this: \textbf{ge-} is unstressed, \textbf{be-} and \textbf{for-} usually are, and a prefixed verb generally stresses its root. In \textbf{gescēop}, “created”, stress \textbf{scēop}, not \textbf{ge-}. A macron indicates duration regardless of which syllable is stressed.

% mdparagraph:P038:end



% mdparagraph:P039:begin

For these careful reading exercises, observe printed length. In connected speech, unstressed function words and prefixes could be shorter than their conventional written macrons suggest. This does not authorize adding, removing, or moving marks in the course text.

% mdparagraph:P039:end



\reserve{5\baselineskip}

% mdparagraph:P040:begin

\section*{7. Practice}

% mdparagraph:P040:end



% mdparagraph:P041:begin

The following tasks are new teaching material. Meanings are supplied throughout. Use the guide's sound convention; no translation or grammatical analysis is required.

% mdparagraph:P041:end



\reserve{5\baselineskip}

% mdparagraph:P042:begin

\subsection*{A. Length contrasts}

% mdparagraph:P042:end



% mdparagraph:P043:begin

Read each pair and identify the long vowel.

% mdparagraph:P043:end



% item: PRON-P01a

% mdparagraph:P044:begin

\textbf{A1.} \textbf{ac}, “but”; \textbf{āc}, “oak”.

% mdparagraph:P044:end



% item: PRON-P01b

% mdparagraph:P045:begin

\textbf{A2.} \textbf{is}, “is”; \textbf{īs}, “ice”.

% mdparagraph:P045:end



% item: PRON-P01c

% mdparagraph:P046:begin

\textbf{A3.} \textbf{god}, “god”; \textbf{gōd}, “good”.

% mdparagraph:P046:end



% item: PRON-P01d

% mdparagraph:P047:begin

\textbf{A4.} \textbf{full}, “full”; \textbf{fūl}, “vile”.

% mdparagraph:P047:end



\reserve{5\baselineskip}

% mdparagraph:P048:begin

\subsection*{B. Consonants and unfamiliar sounds}

% mdparagraph:P048:end



% mdparagraph:P049:begin

Give the taught sound or mouth instruction for the indicated letters, then read the word.

% mdparagraph:P049:end



% item: PRON-P02a

% mdparagraph:P050:begin

\textbf{B1.} \textbf{c} in \textbf{cuman}, “come”.

% mdparagraph:P050:end



% item: PRON-P02b

% mdparagraph:P051:begin

\textbf{B2.} Both \textbf{c} letters in \textbf{cyrice}, “church”.

% mdparagraph:P051:end



% item: PRON-P02c

% mdparagraph:P052:begin

\textbf{B3.} \textbf{c} in \textbf{cēlan}, “cool”.

% mdparagraph:P052:end



% item: PRON-P02d

% mdparagraph:P053:begin

\textbf{B4.} \textbf{g} in \textbf{giefan}, “give”.

% mdparagraph:P053:end



% item: PRON-P02e

% mdparagraph:P054:begin

\textbf{B5.} \textbf{g} in \textbf{fugol}, “bird”.

% mdparagraph:P054:end



% item: PRON-P02f

% mdparagraph:P055:begin

\textbf{B6.} \textbf{sc} in \textbf{scip}, “ship”.

% mdparagraph:P055:end



% item: PRON-P02g

% mdparagraph:P056:begin

\textbf{B7.} \textbf{h} in \textbf{niht}, “night”.

% mdparagraph:P056:end



% item: PRON-P02h

% mdparagraph:P057:begin

\textbf{B8.} \textbf{y} in \textbf{wyrt}, “herb; plant”.

% mdparagraph:P057:end



\reserve{5\baselineskip}

% mdparagraph:P058:begin

\subsection*{C. Voice}

% mdparagraph:P058:end



% mdparagraph:P059:begin

State whether the indicated sound is voiced or unvoiced. Give its English sound cue and explain the context or supplied exception.

% mdparagraph:P059:end



% item: PRON-P03a

% mdparagraph:P060:begin

\textbf{C1.} \textbf{f} in \textbf{cnīf}, “knife”.

% mdparagraph:P060:end



% item: PRON-P03b

% mdparagraph:P061:begin

\textbf{C2.} \textbf{f} in \textbf{cnīfas}, “knives”.

% mdparagraph:P061:end



% item: PRON-P03c

% mdparagraph:P062:begin

\textbf{C3.} \textbf{s} in \textbf{rīsan}, “rise”.

% mdparagraph:P062:end



% item: PRON-P03d

% mdparagraph:P063:begin

\textbf{C4.} \textbf{s} in \textbf{hūs}, “house”.

% mdparagraph:P063:end



% item: PRON-P03e

% mdparagraph:P064:begin

\textbf{C5.} \textbf{ð} in \textbf{ōðer}, “other”.

% mdparagraph:P064:end



% item: PRON-P03f

% mdparagraph:P065:begin

\textbf{C6.} \textbf{ð} in \textbf{hǣlðu}, “health”.

% mdparagraph:P065:end



\reserve{5\baselineskip}

% mdparagraph:P066:begin

\subsection*{D. Syllables and stress}

% mdparagraph:P066:end



% mdparagraph:P067:begin

Give the number of syllables and identify the stressed one. Identify any doubled consonant and explain how it is pronounced.

% mdparagraph:P067:end



% item: PRON-P04a

% mdparagraph:P068:begin

\textbf{D1.} \textbf{habban}, “have”.

% mdparagraph:P068:end



% item: PRON-P04b

% mdparagraph:P069:begin

\textbf{D2.} \textbf{fugol}, “bird”.

% mdparagraph:P069:end



% item: PRON-P04c

% mdparagraph:P070:begin

\textbf{D3.} \textbf{heofonan}, “heaven”.

% mdparagraph:P070:end



% item: PRON-P04d

% mdparagraph:P071:begin

\textbf{D4.} \textbf{gescēop}, “created”.

% mdparagraph:P071:end



\reserve{5\baselineskip}

% mdparagraph:P072:begin

\subsection*{E. A sentence adapted from Cook}

% mdparagraph:P072:end



% mdparagraph:P073:begin

This sentence adapts the opening of Cook's Creation reading.[8]

% mdparagraph:P073:end



% mdparagraph:P074:begin

\textbf{God gescēop heofonan and eorðan.}

% mdparagraph:P074:end



% mdparagraph:P075:begin

“God created heaven and earth.”

% mdparagraph:P075:end



% mdtable:T05:begin
\par\begingroup
\fontsize{10.5}{13.2}\selectfont
\setlength{\LTleft}{0pt}\setlength{\LTright}{\fill}
\setlength{\LTpre}{5pt}\setlength{\LTpost}{5pt}
\begin{longtable}{@{}>{\raggedright\arraybackslash}p{21.000mm}>{\raggedright\arraybackslash}p{24.000mm}>{\raggedright\arraybackslash}p{74.377mm}@{}}
\toprule
\textbf{Word} & \textbf{Supplied meaning} & \textbf{Pronunciation help} \\
\midrule
\endfirsthead
\toprule
\textbf{Word} & \textbf{Supplied meaning} & \textbf{Pronunciation help} \\
\midrule
\endhead
\bottomrule
\endfoot
\bottomrule
\endlastfoot
God & God & One syllable; hard \textbf{g} and short rounded \textbf{o}. \\
gescēop & created & Two syllables; stress the second. Initial \textbf{g} is \textbf{y} as in yes; \textbf{sc} is \textbf{sh}. Keep the first \textbf{e} light, and make long \textbf{ēo} one glide. \\
heofonan & heaven & Three syllables; stress the first. \textbf{eo} is one short glide; \textbf{f} is voiced, as in van. Keep both following vowel syllables. \\
and & and & A light connective; keep the final \textbf{d} audible. \\
eorðan & earth & Two syllables; stress the first. \textbf{eo} is one short glide; \textbf{ð} is voiced, as in this. Keep the final \textbf{-an} audible. \\
\end{longtable}
\par\endgroup
% mdtable:T05:end

% mdparagraph:P076:begin

Read the complete sentence, using the supplied help, then answer these sound questions.

% mdparagraph:P076:end



% item: PRON-P05a

% mdparagraph:P077:begin

\textbf{E1.} Which sounds do \textbf{g} and \textbf{sc} have in \textbf{gescēop}, “created”? Explain its two syllables.

% mdparagraph:P077:end



% item: PRON-P05b

% mdparagraph:P078:begin

\textbf{E2.} Why is \textbf{f} voiced in \textbf{heofonan}, “heaven”?

% mdparagraph:P078:end



% item: PRON-P05c

% mdparagraph:P079:begin

\textbf{E3.} Why is \textbf{ð} voiced in \textbf{eorðan}, “earth”, even though the preceding letter is not a vowel?

% mdparagraph:P079:end



\reserve{5\baselineskip}

% Keep the source notes together as a compact final reference page.
\newpage

% mdparagraph:P080:begin

\section*{Source notes}

% mdparagraph:P080:end

\begingroup
\fontsize{9.5}{12}\selectfont



% mdparagraph:P081:begin

1. Cook, \emph{A First Book in Old English}, second edition (1895), §§2–4, pp.5–6 / PDF21–22; Fulk, \emph{An Introductory Grammar of Old English}, corrected open-access reprint (2020), §§15–16, pp.10–12 / PDF22–24. Fulk's teaching dots are omitted in the selected examples; \textbf{hǣlðu, forðām} replace his thorn with the course's eth. These declared spelling changes do not change their selected sounds.

% mdparagraph:P081:end



% mdparagraph:P082:begin

2. Fulk §§12,16,19, pp.10–11,13–14 / PDF22–23,25–26, permits the traditional short-vowel approximations used here. Held \emph{cat} for \textbf{ǣ} follows Fulk. Cook §5, pp.6–7 / PDF22–23, also offers practical approximations, but some of his English cues suit other accents. American cue adjustments are teaching aids.

% mdparagraph:P082:end



% mdparagraph:P083:begin

3. Fulk §§20–21, p.14 / PDF26; \textbf{giefan}, §135a, p.103 / PDF115. The rounded diphthongs follow Fulk's declared convention. Cook's §5 table, p.7 / PDF23, gives different component vowels; the two descriptions are not interchangeable.

% mdparagraph:P083:end



% mdparagraph:P084:begin

4. Fulk §§16,18,135a–d, pp.11–13,103–104 / PDF23–25,115–116; Cook §§7,10, pp.8–9 / PDF24–25. The named exceptions bound the general spelling tendencies.

% mdparagraph:P084:end



% mdparagraph:P085:begin

5. Cook §§7,10, pp.8–9 / PDF24–25; Fulk §§9,16,18, pp.8,11–13 / PDF20,23–25. American \textbf{r} is a practical approximation. The \textbf{k+s} reading of \textbf{x/hs} follows Cook's practical value; Fulk also records an earlier h-like sound before \textbf{s}. Neither source requires a trill in these practical descriptions.

% mdparagraph:P085:end



% mdparagraph:P086:begin

6. Fulk §17, p.12 / PDF24, supplies the voicing rule and exceptions used here. Cook §8, pp.8–9 / PDF24–25, gives a broader syllable-final voicing rule. This guide selects Fulk's account.

% mdparagraph:P086:end



% mdparagraph:P087:begin

7. Cook §11, pp.9–10 / PDF25–26, distinguishes prolonged medial doubling from final double spellings, normally read with single-consonant duration. See also Cook §23, pp.18–19 / PDF34–35, and Fulk §§17,19,22–23, pp.12–15 / PDF24–27. Printed vowel length and stress are distinct; connected-speech shortening does not require changing the edited spelling.

% mdparagraph:P087:end



% mdparagraph:P088:begin

8. Adapted from Cook, Reader I, p.123 / PDF139, first prose line: remove \textbf{On anginne}, move \textbf{God} before the verb, and omit the three apparatus cues. The retained words are not respelled. The original sentence and notes remain in the source transcription. The adaptation's full meaning and pronunciation help are supplied here.

% mdparagraph:P088:end



\endgroup
\end{document}
